\documentclass[12pt,letterpaper]{article}
\usepackage[margin=0.5in]{geometry}
\usepackage{array}
\usepackage{amsmath}
\usepackage{graphicx}
\renewcommand{\familydefault}{\sfdefault}
\setlength{\parindent}{0pt}
\setlength{\parskip}{0.35em}
\setlength{\fboxsep}{3pt}

\newcolumntype{C}[1]{>{\centering\arraybackslash}m{#1}}
\newcommand{\answerbox}[1]{%
  \begin{center}
    {\bfseries #1}\\
    \vspace*{\fill}
    {\Large Write your answer\\ in your Class Notes}
  \end{center}
  \vspace*{\fill}
}

\begin{document}

\begin{center}
    {\Large\bfseries Assignment 5: Ratio, Proportion, and Percent}\\[4pt]
    {\large\bfseries DO NOT WRITE ON THIS PAPER}\\[4pt]
\end{center}
\vspace{-0.8em}

\noindent\hfill\fbox{%
  \begin{minipage}[t][0.34\textheight][t]{.95\textwidth}
    \begin{center}
      \textbf{Part A}
    \end{center}

    On your whiteboard, draw the table below.

    Flip 2 cards and record them below. Make a ratio of card 1 to card 2. Identify two ratios that are
    proportional to your original ratio. \emph{Hint: How did we double a recipe to make a proportion?}

    \begin{center}
    \renewcommand{\arraystretch}{1.25}
    \begin{tabular}{|C{0.08\textwidth}|C{0.08\textwidth}|C{0.23\textwidth}|C{0.16\textwidth}|C{0.19\textwidth}|}
      \hline
      Card 1 & Card 2 & Ratio of Card 1:Card 2 & First proportion & Second proportion\\
      \hline
      3 & 7 & 3:7 & 6:14 & 9:21\\
      \hline
      \rule{0pt}{1.55em} & & & &\\
      \hline
      \rule{0pt}{1.55em} & & & &\\
      \hline
    \end{tabular}
    \end{center}

    Do one example in front of the TA and explain your steps. On your paper write one example.
    \vfill
  \end{minipage}%
}\hfill

\noindent\hfill\fbox{%
  \begin{minipage}[t][0.48\textheight][t]{.95\textwidth}
    \begin{center}
      \textbf{Part B}
    \end{center}

    Flip 3 cards and record them on your whiteboard.

    \vspace{1.1em}
    \begin{center}
      Card 1 = \hspace{3em} Card 2 = \hspace{3em} Card 3 =
    \end{center}
    \vspace{1.1em}

    Use those cards to set up the following proportions:
    \[
      \frac{x}{\text{card 1}}=\frac{\text{card 2}}{\text{card 3}}
      \qquad
      \frac{\text{card 1}}{x}=\frac{\text{card 2}}{\text{card 3}}
      \qquad
      \frac{\text{card 1}}{\text{card 2}}=\frac{\text{card 3}}{x}
    \]

    Solve each proportion and round your answer to the nearest hundredth.

    Do one example in front of the TA and explain your steps. On your paper write one example of each
    proportion.
    \vfill
  \end{minipage}%
}\hfill

\newpage

\noindent\hfill\fbox{%
  \begin{minipage}[t][0.40\textheight][t]{.95\textwidth}
    \begin{center}
      \textbf{Part C}
    \end{center}

    Flip two cards. Use the first card as your numerator and your second card as your denominator.
    Convert each fraction into a decimal and round to the nearest \textbf{ten-thousandth}. Convert each decimal
    into a percentage. Write each percentage as a ratio out of 100.

    \begin{center}
    \renewcommand{\arraystretch}{1.45}
    \begin{tabular}{|C{0.10\textwidth}|C{0.10\textwidth}|C{0.20\textwidth}|C{0.13\textwidth}|}
      \hline
      \textbf{Fraction} & \textbf{Decimal} & \textbf{Ratio with 100} & \textbf{Percentage}\\
      \hline
      \rule{0pt}{2.15em}$\dfrac{4}{7}$ & 0.5614 & $\dfrac{56.14}{100}$ or 56.14:100 & 56.14\%\\
      \hline
      \rule{0pt}{2.15em}$\dfrac{9}{4}$ & 2.2500 & $\dfrac{225}{100}$ or 225:100 & 225.00\%\\
      \hline
      \rule{0pt}{1.35em} & & &\\
      \hline
      \rule{0pt}{1.35em} & & &\\
      \hline
    \end{tabular}
    \end{center}
    \vfill
  \end{minipage}%
}\hfill

\noindent\hfill\fbox{%
  \begin{minipage}[t][0.46\textheight][t]{.95\textwidth}
    \begin{center}
      \textbf{Question 1}
    \end{center}

    \textbf{Big Question:} What is a percentage? Descibe what a percentage is in terms of a ratio.

    \vspace{0.6em}

    Using one of the problems you solved in part C, describe in words the steps to convert a fraction to a
    percentage using proportions.

    \vspace{0.6em}

    Word bank: \emph{ratio, proportion, equivalent, percentage}
    \vfill
  \end{minipage}%
}\hfill

\newpage

\noindent\hfill\fbox{%
  \begin{minipage}[t][0.42\textheight][t]{.70\textwidth}
    \begin{center}
      \textbf{Question 2 Work}
    \end{center}

    A machine produces 312 bolts in 30 minutes. At the same rate, how many
    minutes would it take to produce 260 bolts?

    Write your answer using a complete sentence.
    \vfill
  \end{minipage}%
}%
\fbox{%
  \begin{minipage}[t][0.42\textheight][t]{0.24\textwidth}
    \answerbox{Question 2 Answer}
  \end{minipage}%
}\hfill

\noindent\hfill\fbox{%
  \begin{minipage}[t][0.42\textheight][t]{.70\textwidth}
    \begin{center}
      \textbf{Question 3 Work}
    \end{center}

    A bank is $\dfrac{3}{6}$ full. After adding 120 pennies, the bank is now $\dfrac{5}{6}$ full. How many
    pennies can the bank hold?

    Draw a tape diagram to represent the bank before and after pennies. Then
    use the diagram to help you solve the problem.

    Write your answer using a complete sentence.
    \vfill
  \end{minipage}%
}%
\fbox{%
  \begin{minipage}[t][0.42\textheight][t]{0.24\textwidth}
    \answerbox{Question 3 Answer}
  \end{minipage}%
}\hfill

\newpage

\begin{center}
    \textbf{Additional Space}
\end{center}

\vspace{0.4em}

\noindent\makebox[\textwidth][c]{%
\fbox{%
  \rule{0pt}{0.20\textheight}\rule{7.35in}{0pt}%
}%
}

\vspace{0.45em}

\noindent\makebox[\textwidth][c]{%
\fbox{%
  \rule{0pt}{0.20\textheight}\rule{7.35in}{0pt}%
}%
}

\vspace{0.45em}

\noindent\makebox[\textwidth][c]{%
\fbox{%
  \rule{0pt}{0.20\textheight}\rule{7.35in}{0pt}%
}%
}

\vspace{0.45em}

\noindent\makebox[\textwidth][c]{%
\fbox{%
  \rule{0pt}{0.20\textheight}\rule{7.35in}{0pt}%
}%
}

\end{document}
